\section{偶数冪根の二頂点角とMcKay構造}
\label{sec:notes-root-McKay-v38}

偶数位数 $\ell=2d$ のchartでは、中心の有限エタール拡大
$\lambda^{2d}=\Xi$ を取ると、二頂点角が全体をMorita的に代表する。
$x$ の固有冪等元を $e_i$ とし、
\[
 e=e_0+e_d,\quad g=e_0-e_d,\quad P=eu^de,\quad R=ev^de
\]
と置けば
\[
 g^2=e,\quad gP=-Pg,\quad gR=-Rg,\quad
 PR=RP=f=(-1)^dQ^d+\beta^d\lambda^{2d}.
\]
一般化Weyl代数の次数が $d\mathbb Z$ のときだけこの角に残るので、
これらの関係は完全である。従って
\[
 eA'e\simeq
 \left(\frac{\mathbb C[Q,J,\lambda^{\pm1},P,R]}{(PR-f)}\right)\#\mu_2.
\]
括弧内の可換環は $\partial_\lambda f$ が単元なので滑らかである。
$e_0,e_d$ は、特異ファイバーの二本の $d$ 頂点鎖を一つずつ代表し、
滑らかな中心ファイバーでもfullである。
有限中心上の全閉点でfullであることから、$e$ は大域的にもfullとなる。
従って有限エタール基底変換の後で、chart全体が滑らかな二元群商スタックの
加群圏で記述できる。元のchartの大域的な行列環表示は、この議論には含まない。

特異中心の完成では $Q\ne0$ なので $PR=f$ から $Q$ を消去できる。
二本の鎖の可逆辺を縮約すると
\[
 \widehat A_p\simeq
 M_d\bigl(\mathbb C[[a,b,z_1,z_2]]\#\mu_2\bigr),
 \qquad (a,b)\longmapsto(-a,-b)
\]
を得る。中心は
$\mathbb C[[a^2,ab,b^2,z_1,z_2]]$、横断特異点は $A_1$ である。
作用は余次元一で自由なので、skew group algebraは
正規不変環上の自己準同型環
$\operatorname{End}_{L^{\mu_2}}(L)$、$L=\mathbb C[[a,b,z_1,z_2]]$ と一致する。
最大Cohen--Macaulay性と大域次元四により、完成代数は非可換クレパント解消となる。
その特殊ファイバーが
$M_d((\mathbb C[a,b]/(a,b)^2)\#\mu_2)$ であり、
二単純加群と平方零の四矢を同時に説明する。

この計算の再利用点は、巡回一般化Weyl代数の係数 $a_i$ が周期 $d$ を持ち、
零辺が二つであるとき、反対の二頂点を選ぶだけで、
長い巡回quiverを二元群のskew group algebraに置き換えられることである。
位数を変えても形式Morita型は変わらず、行列サイズだけが変わる。
